% FILE: common/scholia/tests/test-interactive.tex
% ==================================================================================================
% TEST FILE: Interactive Hide/Reveal Features
%
% PURPOSE:
%   Test the OCG-based hide/reveal functionality provided by interactive/main.sty
%
% REQUIREMENTS:
%   - pdflatex, lualatex, or xelatex
%   - PDF viewer with OCG/JavaScript support (Adobe Reader, PDF.js)
%   - Note: Preview.app on macOS does NOT support OCG toggling
%
% USAGE:
%   pdflatex test-interactive.tex
%   Open test-interactive.pdf in Adobe Reader
%   Click on dotted underlines to reveal hidden content
%
% ==================================================================================================
\documentclass[11pt,a4paper]{article}

% Path to common directory (parent of scholia)
\newcommand{\CommonPath}{../..}

% Load scholia with interactive option
\usepackage[interactive,french]{\CommonPath/scholia/scholia}

% Additional packages for testing
\usepackage{amsmath}
\usepackage{lipsum}

\title{Test: Interactive Hide/Reveal Features}
\author{Scholia Package Tests}
\date{\today}

\begin{document}
\maketitle

\section{Introduction}

This document tests the interactive hide/reveal functionality provided by the
\texttt{scholia} package with the \texttt{interactive} option.

\textbf{Instructions:}
\begin{enumerate}
  \item Open this PDF in a viewer that supports OCG layers (Adobe Reader, PDF.js)
  \item Click on the dotted underlined areas to reveal hidden content
  \item Click again to hide the content
\end{enumerate}

Note: Apple Preview and some other viewers do not support OCG toggling.

% ==================================================================================================
\section{Basic Hide Command}
% ==================================================================================================

\subsection{Simple Text}

The capital of France is \hide{Paris}.

The answer to $2 + 2$ is \hide{4}.

A complete sentence can be hidden: \hide{This entire sentence is hidden until clicked.}

\subsection{Multiple Hidden Elements}

Fill in the blanks:
\begin{itemize}
  \item The largest planet is \hide{Jupiter}.
  \item Water boils at \hide{100} degrees Celsius.
  \item The chemical symbol for gold is \hide{Au}.
  \item Light travels at approximately \hide{300,000 km/s}.
\end{itemize}

% ==================================================================================================
\section{Math Mode Compatibility}
% ==================================================================================================

\subsection{Inline Math}

The quadratic formula is $x = \hide{\frac{-b \pm \sqrt{b^2 - 4ac}}{2a}}$.

Given $f(x) = x^2$, we have $f'(x) = \hide{2x}$.

The derivative of $\sin(x)$ is $\hide{\cos(x)}$.

\subsection{Display Math}

Solve for $x$:
\[
  x^2 - 5x + 6 = 0 \implies x = \hide{2} \text{ or } x = \hide{3}
\]

The integral:
\[
  \int_0^1 x^2 \, dx = \hide{\left[\frac{x^3}{3}\right]_0^1} = \hide{\frac{1}{3}}
\]

\subsection{Aligned Equations}

\begin{align*}
  (a + b)^2 &= \hide{a^2 + 2ab + b^2} \\
  (a - b)^2 &= \hide{a^2 - 2ab + b^2} \\
  (a + b)(a - b) &= \hide{a^2 - b^2}
\end{align*}

% ==================================================================================================
\section{Custom Colors}
% ==================================================================================================

The default reveal color is black. You can specify a custom color:

\begin{itemize}
  \item Default color: \hide{This text appears in black}
  \item Red text: \hide[red]{This text appears in red}
  \item Blue text: \hide[blue]{This text appears in blue}
  \item Green text: \hide[green!60!black]{This text appears in dark green}
  \item Custom color: \hide[orange]{This text appears in orange}
\end{itemize}

% ==================================================================================================
\section{Nested and Complex Content}
% ==================================================================================================

\subsection{Formatted Text}

Hidden content can include formatting:
\begin{itemize}
  \item Bold: \hide{\textbf{bold text}}
  \item Italic: \hide{\textit{italic text}}
  \item Combined: \hide{\textbf{\textit{bold italic}}}
\end{itemize}

\subsection{Multiple Words}

The three primary colors are \hide{red, blue, and yellow}.

In French: Les trois couleurs primaires sont \hide{le rouge, le bleu et le jaune}.

% ==================================================================================================
\section{Integration with Scholia Features}
% ==================================================================================================

\subsection{With Callout Boxes}

\begin{definitionbox}{Pythagorean Theorem}
  In a right triangle with legs $a$ and $b$ and hypotenuse $c$:
  \[
    \hide{a^2 + b^2 = c^2}
  \]
\end{definitionbox}

\begin{methodbox}{Solving Quadratic Equations}
  To solve $ax^2 + bx + c = 0$:
  \begin{enumerate}
    \item Calculate the discriminant: $\Delta = \hide{b^2 - 4ac}$
    \item If $\Delta > 0$: two solutions $x = \hide{\frac{-b \pm \sqrt{\Delta}}{2a}}$
    \item If $\Delta = 0$: one solution $x = \hide{\frac{-b}{2a}}$
    \item If $\Delta < 0$: \hide{no real solutions}
  \end{enumerate}
\end{methodbox}

\subsection{With Exercise Environment}

\begin{exauto}[Derivatives][1]
  Calculate the following derivatives:
  \begin{enumerate}
    \item $f(x) = 3x^4 - 2x^2 + 5$ \quad $\Rightarrow$ \quad $f'(x) = \hide{12x^3 - 4x}$
    \item $g(x) = \sin(2x)$ \quad $\Rightarrow$ \quad $g'(x) = \hide{2\cos(2x)}$
    \item $h(x) = e^{x^2}$ \quad $\Rightarrow$ \quad $h'(x) = \hide{2x \cdot e^{x^2}}$
  \end{enumerate}
\end{exauto}

% ==================================================================================================
\section{Edge Cases}
% ==================================================================================================

\subsection{Empty Content}

An empty hide command: \hide{} (should produce nothing).

\subsection{Special Characters}

Hide with special characters: \hide{10\% discount, \$100, \& more}.

\subsection{Long Content}

A longer hidden passage: \hide{This is a much longer piece of text that demonstrates
how the hide command handles content that spans multiple words and could potentially
wrap across lines in the final document.}

\subsection{Adjacent Hidden Elements}

Two hidden elements side by side: \hide{first}\hide{second}.

With space: \hide{first} \hide{second}.

% ==================================================================================================
\section{Performance Check}
% ==================================================================================================

This section tests multiple hidden elements to ensure the counter mechanism works:

\begin{enumerate}
  \item Answer: \hide{A}
  \item Answer: \hide{B}
  \item Answer: \hide{C}
  \item Answer: \hide{D}
  \item Answer: \hide{E}
  \item Answer: \hide{F}
  \item Answer: \hide{G}
  \item Answer: \hide{H}
  \item Answer: \hide{I}
  \item Answer: \hide{J}
\end{enumerate}

% ==================================================================================================
\section{Conclusion}
% ==================================================================================================

All tests above should demonstrate:
\begin{itemize}
  \item Basic hide/reveal functionality with clickable toggling
  \item Proper handling of text and math modes
  \item Custom color support
  \item Integration with other scholia features
  \item Correct counter management for unique layer IDs
\end{itemize}

If all hidden content can be revealed by clicking, the test is successful.

\end{document}
